% 五套模板共用同一正文骨架；换版式时保留本文件及全部实测材料。
\section{引言}
\ReportTODO{约200--400字，说明背景、目的、意义与研究问题；引用可靠来源，如讲义\cite{manual}（此处为占位）。}

\section{原理}
\ReportTODO{约300--800字，自行归纳模型、假设、关系与适用范围，区分物理、实验和仪器原理。}
公式须编号并定义符号。例如，模型 $y=f(x_1,\ldots,x_n)$ 在线性近似且输入相互独立时，合成标准不确定度为
\begin{equation}
  u_c^2(y)=\sum_{i=1}^{n}
    \left(\frac{\partial f}{\partial x_i}\right)^2 u^2(x_i).
  \label{eq:uncertainty}
\end{equation}
式\eqref{eq:uncertainty}为示例。\ReportTODO{补充实际模型、符号和单位；输入相关时加协方差项，不适用时替换。}

\section{实验}
\subsection{仪器与实验条件}
\ReportTODO{列出实际仪器型号、量程、分辨率、接法、环境和样品状态；未知信息不估造。}
\subsection{方法与实际步骤}
\ReportTODO{按实际顺序写测量方法、设置、重复次数、读数与现象，记录偏离讲义的操作及异常。}

\section{结果与分析讨论}
\subsection{定量结果与原始记录}
\ReportTODO{区分原始、计算、理论和估计值，给单位、有效数字及代表性计算。表\ref{tab:raw}仅为格式，无实测数据。}
\begin{table}[H]
  \centering\small
  \caption{原始记录占位；按实测替换列名及单位，保留缺测说明。}
  \label{tab:raw}
  \begin{tabularx}{\columnwidth}{@{}cXXX@{}}
    \toprule
    序号 & 物理量/单位 & 实测读数 & 状态说明\\
    \midrule
    --- & 待补充 & 待补充 & 无原始数据\\
    \bottomrule
  \end{tabularx}
\end{table}
\subsection{规律与进一步处理}
\ReportTODO{说明定标、拟合、统计或插值；给模型、参数、单位及R\textsuperscript{2}或残差。图\ref{fig:result}须由真实数据生成。}
\begin{figure}[H]
  \centering
  \fbox{\parbox[c][22mm][c]{0.9\columnwidth}{\centering 待补充：真实数据图\\此框仅为版面占位}}
  \caption{结果图占位。坐标轴、图内字体与图题须标清物理量和单位。}
  \label{fig:result}
\end{figure}
\subsection{与理论比较并解释}
\ReportTODO{与有来源的理论或标准比较，给偏差和适用条件，解释实测现象。}
\subsection{误差来源与改进建议}
\ReportTODO{分析系统、随机、读数、环境和仪器误差。保留异常值，剔除须给判据与理由；至少两条对应误差来源的可执行改进。}

\section{结论和建议（总结）}
\ReportTODO{回答实验目的，归纳有证据的结果及规律，给单位和适当不确定度，说明局限及建议，不引入未经论证的新结果。}

% 文献表须只保留正文实际引用的3--6条资料（教师另有规定时以其为准）。
% 占位条目 manual 不是已核实文献，正式提交前必须替换。
\renewcommand{\refname}{六、参考文献}
\bibliographystyle{gbt7714-numerical}
\bibliography{report_template/references}

\ifReportAppendix
\ifdefined\ReportEndColumns\ReportEndColumns\fi
\section*{附录：原始记录与复现信息}
\ReportTODO{按需附原始记录、数据索引、程序、计算依据和补充图表。材料分别在data/、scripts/、figures/；无材料可删附录。}
\fi
